\section*{Chapter \ref{ch:s1:options-implied-signals-for-stocks} --- Options-Implied Signals for Stocks}

\begin{solution}{exo:s1:options-implied-signals-for-stocks:1}
Bad news is weighted 1.5: $0.003 \times 1.5 \times 2 = 0.009$, 0.9 volatility points on the puts; good news $0.003 \times 2 = 0.006$, 0.6 points on the calls. Both
are below the daily noise of one point, which is why single days are uninformative and the signal must be read against its history.
\end{solution}

\begin{solution}{exo:s1:options-implied-signals-for-stocks:2}
$0.34\% \times 52 = 17.7\%$ a year and $0.5\% \times 52 = 26\%$ a year, before costs and without compounding.
\end{solution}

\begin{solution}{exo:s1:options-implied-signals-for-stocks:3}
A steep skew means puts are bid up, as they are when informed traders expect a fall: the stock is more likely to fall, so the skew enters with a minus. A high
call--put spread means calls are bid up relative to puts, as they are when good news is expected: it enters with a plus.
\end{solution}

\begin{solution}{exo:s1:options-implied-signals-for-stocks:4}
Shorting the stock costs a borrow fee, needs a locate, and may be impossible or recalled; buying a put has none of these costs and adds leverage. For good
news, buying the stock is cheap and easy, so options offer only leverage.
\end{solution}

\begin{solution}{exo:s1:options-implied-signals-for-stocks:5}
The informed traders act only before events, so the signals contain information only there; elsewhere they are noise that dilutes the IC. The book should
trade the signals in stocks with an event ahead and stay out of the rest.
\end{solution}

\begin{solution}{exo:s1:options-implied-signals-for-stocks:6}
The implied volatilities carry daily noise of a point, so the weekly ranking changes a lot from week to week: the book trades heavily for a small, diffuse
signal, and ten basis points per unit traded costs more than the 8.5\% a year it earns gross.
\end{solution}

\begin{solution}{exo:s1:options-implied-signals-for-stocks:7}
The volume book around events earns 9.1\% a year after costs (Sharpe ratio 1.75) against 23.0\% and 31.9\% for skew and spread. Volume rises with the size of
the informed demand whatever its direction; only the extra weight on bad news gives it a sign, so its IC is about half that of the price signals.
\end{solution}

\begin{solution}{exo:s1:options-implied-signals-for-stocks:8}
The number comes from a model in which the informed traders know the whole surprise and their demand is a clean function of it. Real informed flow is small,
unsigned in most data, mixed with hedging, and competed for; the published effects are tens of basis points a week before costs. Expect a small fraction,
and test on real option data with costs.
\end{solution}

\begin{solution}{pb:s1:options-implied-signals-for-stocks:1}
\begin{enumerate}
\item Option trading by the privately informed; call minus put implied volatility at matched strikes; put volume over call volume; option volume over stock
volume.
\item Steepest smirks underperform flattest by 10.9\% a year risk-adjusted, for at least six months, with worse earnings shocks next quarter.
\item Expensive calls beat expensive puts by 50 basis points a week, declining over time; the call--put spread predicts positively, realised minus implied
negatively.
\item Low \omterm{def:s1:options-implied-signals-for-stocks:volume}{put--call ratios} beat high by 40 basis points a day and over 1\% a week; low option-to-stock volume beats high by 0.34\% a week.
\item At-the-money volatility equal to realised plus two points, a three-point put skew, a point of daily noise; informed traders buying puts or calls in the
ten days before announcements, 0.3 points per standard deviation of surprise.
\item Short-sale costs make options more attractive for bad news.
\item It carries no planted information, so its IC should be zero, and is.
\item To mark the ten days before each announcement, when the informed traders act.
\item Skew 0.004, 0.014, 0.018, 0.012; spread 0.005, 0.017, 0.022, 0.015; volume 0.002, 0.008, 0.010, 0.007; implied minus realised zero.
\item Skew 0.026, 0.085, 0.111, 0.082; spread 0.030, 0.101, 0.129, 0.094; volume 0.015, 0.047, 0.061, 0.045.
\item The announcement's jump falls within twenty days of the informed buying; beyond, the window adds noise.
\item All stocks: gross Sharpe ratios of 3.17--3.67, negative after costs; around events: 23.0\%, 31.9\% and 9.1\% a year after costs for skew, spread and volume.
\item \textbf{Named result.} The volatility spread's IC is 0.005, 0.017, 0.022 and 0.015 at 1, 5, 20 and 60 days across all stocks and 0.030, 0.101, 0.129 and
0.094 among stocks announcing within ten days; the skew's 0.004, 0.014, 0.018, 0.012 and 0.026, 0.085, 0.111, 0.082.
\item The model's informed traders know the surprise exactly and are the only source of informed demand.
\item Make informed flow a small share of option volume, add hedging and structured flow, and let informed traders choose the stock when it is cheaper.
\item Smooth the signals, trade them only where events are near, and rebalance less often.
\item Implied volatility surfaces by stock, option volume (signed if possible), an event calendar, and borrow costs.
\item The call--put spread (Cremers and Weinbaum).
\item They trade the same events: the options signals anticipate the surprise that chapter~7's drift book trades after it.
\item Because informed traders often trade options first.
\end{enumerate}
\end{solution}

\begin{solution}{iq:s1:options-implied-signals-for-stocks:1}
Informed traders use options for leverage and to avoid short-sale costs; their demand moves option prices and volumes before the stock moves, and the stock
market is slow to follow.
\end{solution}

\begin{solution}{iq:s1:options-implied-signals-for-stocks:2}
From a fitted surface at a fixed maturity (for example 30 days): the implied volatility at a fixed out-of-the-money moneyness (a delta of $-0.2$ or a strike at
90\%) minus the at-the-money implied volatility, from closing quotes, compared with the stock's own history and the market's skew.
\end{solution}

\begin{solution}{iq:s1:options-implied-signals-for-stocks:3}
The call--put spread compares two options on the same stock and measures relative demand for upside against downside, an information signal; implied minus
realised compares options with the stock's own history and measures the volatility risk premium, a risk signal.
\end{solution}

\begin{solution}{iq:s1:options-implied-signals-for-stocks:4}
Check whether it is buyer-initiated and opening, compare it with the usual pre-earnings volume, the skew and the stock's short interest; if it looks like
informed buying of protection, reduce or hedge longs, and consider the stock underweight into the event.
\end{solution}

\begin{solution}{iq:s1:options-implied-signals-for-stocks:5}
Stale or non-synchronous quotes between options and stock; bid--ask bounce in implied volatilities; early-exercise premia; survivorship in option data; and
costs of trading the stock on noisy signals.
\end{solution}

\begin{solution}{iq:s1:options-implied-signals-for-stocks:6}
With private information of a fall of size $x$, shorting earns $x - c$ per unit; a put with leverage $L$ earns roughly $Lx$ minus its premium's time value. The
informed prefer puts when $Lx$ net of premium exceeds $x - c$, which is more likely the higher $c$: when shorting is costly, bad news moves into puts, so put
volume (and the \omterm{def:s1:options-implied-signals-for-stocks:volume}{option-to-stock volume ratio}) carries more negative information than call volume carries positive.
\end{solution}
